\notechapter{N-20221012}{A Map of Classical Inequalities}

\section{Bernoulli's inequality}

The note begins in the middle of a list of standard inequalities.  Bernoulli's inequality says that if \(x>-1\) and \(n\in\mathbb N\), then
\[
  (1+x)^n\ge 1+nx.
\]
A slightly more general form is: if \(x_i>-1\) and the \(x_i\)'s have the same sign, then
\[
  (1+x_1)(1+x_2)\cdots(1+x_n)
  \ge 1+x_1+x_2+\cdots+x_n.
\]
When \(x_1=\cdots=x_n=x\), this reduces to Bernoulli's inequality.

The handwritten note adds the useful memory: for \(x>0\), the graph of \(y=(1+x)^n\) lies above its tangent line at \(0\).  This is the convexity explanation.

\section{Power means}

For a positive vector
\[
  a=(a_1,\ldots,a_n),
\]
and \(r\neq0\), define the power mean
\[
  M_r(a)=\left(\frac{a_1^r+a_2^r+\cdots+a_n^r}{n}\right)^{1/r}.
\]
The power mean inequality says that if \(r\le s\), then
\[
  M_r(a)\le M_s(a).
\]
Thus the familiar chain is
\[
  \text{harmonic mean}\le\text{geometric mean}\le\text{arithmetic mean}\le\text{quadratic mean}.
\]

This is a good example of a high-level unification.  Instead of remembering many separate inequalities, one remembers that changing the exponent changes the mean monotonically.

\section{Cauchy--Schwarz}

The Cauchy inequality is
\[
  (a_1^2+\cdots+a_n^2)(b_1^2+\cdots+b_n^2)
  \ge (a_1b_1+\cdots+a_nb_n)^2.
\]
It is equality exactly when the two vectors are proportional:
\[
  (a_1,\ldots,a_n)=\lambda(b_1,\ldots,b_n).
\]

The Engel form, also called Titu Andreescu's lemma, is
\[
  \frac{a_1^2}{b_1}+\frac{a_2^2}{b_2}+\cdots+\frac{a_n^2}{b_n}
  \ge
  \frac{(a_1+a_2+\cdots+a_n)^2}{b_1+b_2+\cdots+b_n},
\]
for \(b_i>0\).  The note marks this formula because it is one of the most frequently used forms in contest inequalities.

The equality condition is
\[
  \frac{a_1}{b_1}=\frac{a_2}{b_2}=\cdots=\frac{a_n}{b_n}.
\]

\section{Minkowski inequality}

The page also records the inequality
\[
  \sqrt{a_1^2+b_1^2}+\cdots+\sqrt{a_n^2+b_n^2}
  \ge
  \sqrt{(a_1+\cdots+a_n)^2+(b_1+\cdots+b_n)^2}.
\]
This is the triangle inequality in \(\mathbb R^2\), applied repeatedly.  Each pair \((a_i,b_i)\) is a vector, and the left side is the sum of lengths, while the right side is the length of the sum.

This viewpoint is much better than memorizing the expression.  Whenever square roots of sums of squares appear, try to see vectors.

\section{A three-variable inequality}

One of the listed consequences has the form
\[
  \frac{x}{y+z}(a+b)+\frac{y}{z+x}(c+a)+\frac{z}{x+y}(a+b)
  \ge \sqrt{3(ab+bc+ca)},
\]
under positivity assumptions.  The exact printed expression is somewhat compressed, but the intended family is clear: combine Cauchy/Engel form with AM--GM to convert fractions into symmetric expressions.

The practical rule is: denominators such as \(x+y\), \(y+z\), \(z+x\) usually suggest either Nesbitt-type transformations or Engel form.

\section{Chebyshev's inequality}

If
\[
  a_1\le a_2\le\cdots\le a_n,\qquad
  b_1\le b_2\le\cdots\le b_n,
\]
then Chebyshev's sum inequality says
\[
  \left(\frac{a_1+\cdots+a_n}{n}\right)
  \left(\frac{b_1+\cdots+b_n}{n}\right)
  \le
  \frac{a_1b_1+\cdots+a_nb_n}{n}.
\]
Equivalently,
\[
  (a_1+\cdots+a_n)(b_1+\cdots+b_n)
  \le n(a_1b_1+\cdots+a_nb_n).
\]
The note writes this as ``same monotonic order gives larger average product.''  This is the right intuition: sorted sequences positively correlate.

\section{Jensen's inequality}

For a convex function \(f\), Jensen's inequality says
\[
  f(\alpha x+(1-\alpha)y)
  \le \alpha f(x)+(1-\alpha)f(y),
\]
where \(0\le\alpha\le1\).  More generally, if \(\alpha_i\ge0\) and \(\sum\alpha_i=1\), then
\[
  f(\alpha_1x_1+\cdots+\alpha_nx_n)
  \le
  \alpha_1f(x_1)+\cdots+\alpha_nf(x_n).
\]
The note also records the second-derivative criterion: if
\[
  f''(x)\ge0,
\]
then \(f\) is convex.

Jensen is not just another inequality; it is a language.  AM--GM follows by applying Jensen to \(\ln x\), and many trigonometric inequalities follow by checking convexity or concavity on the relevant interval.

\section{Young, Hölder, and Minkowski}

Young's inequality says that if \(a,b>0\), \(p,q>1\), and
\[
  \frac1p+\frac1q=1,
\]
then
\[
  ab\le\frac{a^p}{p}+\frac{b^q}{q}.
\]
Equality holds when
\[
  a^p=b^q.
\]

Hölder's inequality says
\[
  a_1b_1+\cdots+a_nb_n
  \le
  \left(a_1^p+\cdots+a_n^p\right)^{1/p}
  \left(b_1^q+\cdots+b_n^q\right)^{1/q}.
\]
Minkowski's inequality says, for \(p>1\),
\[
  \left(\sum_{i=1}^n(a_i+b_i)^p\right)^{1/p}
  \le
  \left(\sum_{i=1}^n a_i^p\right)^{1/p}
  +
  \left(\sum_{i=1}^n b_i^p\right)^{1/p}.
\]

The note lists these at the end because they are higher-level inequalities.  They belong naturally to normed vector spaces and \(L^p\) spaces.


\section{Cauchy--Schwarz as a master inequality}

The Cauchy--Schwarz inequality
\[
  \left(\sum_i a_i b_i\right)^2\le\left(\sum_i a_i^2\right)\left(\sum_i b_i^2\right)
\]
can be read geometrically as
\[
  |a\cdot b|\le \|a\|\,\|b\|.
\]
Many inequalities are disguised Cauchy--Schwarz.  The skill is to choose the two vectors correctly.

\section{Equality conditions}

For every inequality, the equality condition is part of the theorem.  In Cauchy--Schwarz, equality holds iff the two vectors are linearly dependent.  In AM--GM, equality holds iff all the compared positive numbers are equal.  Tracking equality conditions is essential in optimization problems.

\section{The lesson behind it}

This page is best read as a map rather than a list.  Bernoulli comes from convexity of powers.  Power means organize all classical means.  Cauchy is inner-product geometry.  Minkowski is the triangle inequality.  Chebyshev is monotone correlation.  Jensen is convexity.  Young and Hölder are the algebra of dual exponents.  Once one sees these meanings, choosing the right inequality becomes much less arbitrary.

\section{The hierarchy of means}

Power means form a scale:
\[
  M_r(x)=\left(\frac1n\sum_i x_i^r\right)^{1/r}.
\]
For positive variables, \(M_r\le M_s\) when \(r<s\).  The arithmetic, geometric, and harmonic means are special points on this scale.

This hierarchy organizes many inequalities: instead of proving each comparison separately, one proves monotonicity of the mean parameter.

\section{Convex order and Jensen}

Jensen's inequality is the central convexity principle:
\[
  f(\mathbb E X)\le \mathbb E f(X)
\]
for convex \(f\).  In finite form this is the weighted inequality.  It tells us that spreading a variable increases the expectation of a convex function.

This idea connects Jensen, Karamata, and variance-type estimates.

\section{Minkowski and norm geometry}

Minkowski's inequality is the triangle inequality in \(L^p\):
\[
  \left(\sum_i |x_i+y_i|^p\right)^{1/p}\le \left(\sum_i |x_i|^p\right)^{1/p}+\left(\sum_i |y_i|^p\right)^{1/p}.
\]
It is not merely an algebraic estimate; it says \(\ell^p\) is a normed vector space.

\section{Classical inequalities in one convex landscape}

Classical inequalities are best organized by structure: Bernoulli by convexity of powers, Cauchy and Holder by dual norms, Minkowski by norm geometry, Jensen by convex expectation, and Karamata by majorization.  The named inequalities are landmarks in one landscape.
